%
% teil1.tex -- Komplexe Zahlen als Punkte der Ebene
%
\section{Komplexe Zahlen als Punkte der Ebene
\label{komplex:section:zahlenebene}}
\kopfrechts{Komplexe Zahlen als Punkte der Ebene}
Der gesamte folgende Aufbau ruht auf einer einzigen Identifikation:
Eine komplexe Zahl wird gleichzeitig als algebraisches Objekt
mit Real- und Imaginärteil und als Punkt einer Ebene gelesen.
Beide Lesarten werden im Folgenden parallel geführt,
sodass jede Rechnung eine geometrische Bedeutung erhält.
Die Unterabschnitte stellen zunächst die kartesische Form
\eqref{komplex:eq:kartesisch} bereit
und führen anschliessend Addition, Multiplikation, Betrag und Konjugation
ein.

%
% Algebraische Form und gausssche Zahlenebene
%
\subsection{Algebraische Form und gausssche Zahlenebene
\label{komplex:subsection:algebraisch}}
Eine komplexe Zahl ist ein Ausdruck der Form
\begin{equation}
z
=
a + bi
\label{komplex:eq:kartesisch}
\end{equation}
mit reellen Zahlen $a$ und $b$.
Die Grösse $i$ erfüllt definitionsgemäss $i^2 = -1$
und wird als \emph{imaginäre Einheit} bezeichnet.
\index{imaginare Einheit@imaginäre Einheit}%
Die reelle Zahl $a$ heisst \emph{Realteil} und wird mit $\Re(z)$ bezeichnet,
\index{Realteil}%
$b$ heisst \emph{Imaginärteil} und wird mit $\Im(z)$ bezeichnet.
\index{Imaginarteil@Imaginärteil}
Die Menge aller komplexen Zahlen wird mit $\mathbb{C}$ bezeichnet.

Identifiziert man $z = a + bi$ mit dem Punkt $(a, b)$ der reellen Ebene,
so ist $\mathbb{C}$ als Menge mit $\mathbb{R}^2$ identisch.
Diese Identifikation heisst \emph{gausssche Zahlenebene}.
\index{Zahlenebene, gausssche}%
Die reelle Achse trägt die Zahlen mit $\Im(z) = 0$,
die imaginäre Achse die Zahlen mit $\Re(z) = 0$.
Was $\mathbb{C}$ von $\mathbb{R}^2$ unterscheidet,
ist nicht die Menge der Punkte,
sondern die zusätzliche Multiplikation,
die jedem Punktepaar einen weiteren Punkt zuordnet.
Diese Multiplikation hat eine geometrische Bedeutung,
die in Abschnitt~\ref{komplex:subsection:multiplikation} sichtbar wird.

%
% Addition, Subtraktion und Translation
%
\subsection{Addition, Subtraktion und Translation
\label{komplex:subsection:addition}}
Sind $z_1 = a_1 + b_1 i$ und $z_2 = a_2 + b_2 i$ zwei komplexe Zahlen,
so ist die Summe komponentenweise definiert:
\begin{equation}
z_1 + z_2
=
(a_1 + a_2) + (b_1 + b_2) \, i .
\label{komplex:eq:addition}
\end{equation}
Geometrisch entspricht die Addition genau der Vektoraddition in
$\mathbb{R}^2$.
Stellt man $z_1$ und $z_2$ als Pfeile vom Ursprung zum jeweiligen Punkt dar,
so ergibt sich $z_1 + z_2$ durch Aneinanderfügen dieser Pfeile nach dem
Parallelogrammgesetz.
Insbesondere ist die Abbildung $z \mapsto z + w$ für ein festes $w$
eine \emph{Translation} der Ebene um den Vektor $w$.
\index{Translation}%
Translationen sind die einfachsten Transformationen,
die in dieser Arbeit auftreten.

Die Subtraktion $z_1 - z_2$ hat eine besonders nützliche Bedeutung:
Der Vektor von Punkt $z_2$ zu Punkt $z_1$ wird gerade durch $z_1 - z_2$
dargestellt.
Diese Beobachtung ist klein, aber zentral,
denn sie erlaubt es, Strecken und Richtungen unmittelbar als komplexe
Zahlen zu schreiben.

%
% Multiplikation als Drehstreckung
%
\subsection{Multiplikation als Drehstreckung
\label{komplex:subsection:multiplikation}}
Die Multiplikation ergibt sich aus der Forderung $i^2 = -1$ und der
Distributivität:
\begin{equation}
(a_1 + b_1 i)(a_2 + b_2 i)
=
(a_1 a_2 - b_1 b_2) + (a_1 b_2 + a_2 b_1) \, i .
\label{komplex:eq:multiplikation-kartesisch}
\end{equation}
In dieser Form sieht man der Multiplikation ihre geometrische Bedeutung
kaum an.
Um sie sichtbar zu machen, betrachtet man zwei Spezialfälle.
Multiplikation mit einer reellen Zahl $r > 0$ streckt jeden Punkt vom
Ursprung weg um den Faktor $r$.
Multiplikation mit $i$ bildet $a + bi$ auf $-b + ai$ ab,
was geometrisch einer Drehung um $90^\circ$ gegen den Uhrzeigersinn
entspricht.
Die volle geometrische Deutung der Multiplikation als \emph{Drehstreckung}
\index{Drehstreckung}%
wird in Abschnitt~\ref{komplex:subsection:rotation} abgeleitet.

%
% Betrag, Konjugation und Abstand
%
\subsection{Betrag, Konjugation und Abstand
\label{komplex:subsection:betrag}}
Der \emph{Betrag} einer komplexen Zahl ist definiert als
\begin{equation}
|z|
=
\!\sqrt{a^2 + b^2}
\label{komplex:eq:betrag}
\end{equation}
und entspricht dem euklidischen Abstand des Punktes $z$ vom Ursprung.
Die \emph{konjugiert komplexe} Zahl ist $\overline{z} = a - bi$;
\index{konjugiert komplex}%
sie ergibt sich durch Spiegelung an der reellen Achse.
Aus der Definition folgt unmittelbar
\begin{equation}
z \cdot \overline{z}
=
a^2 + b^2
=
|z|^2 ,
\label{komplex:eq:konjugation}
\end{equation}
was später in Beweisen wiederholt verwendet wird.

Der Abstand zweier Punkte $z_1$ und $z_2$ der gaussschen Ebene ist gerade
$|z_1 - z_2|$.
Damit lassen sich geometrische Bedingungen wie ``der Punkt $z$ hat von zwei
festen Punkten denselben Abstand'' sofort als algebraische Gleichungen
schreiben --- eine Idee, die in Abschnitt~\ref{komplex:section:objekte}
systematisch genutzt wird.
